\section{卷三 73-89}
\subsection{第73篇(亞薩)：在聖所看到惡人結局而不再嫉妒，并敬畏神  73:1-28}
\subsubsection{亚萨因惡人发达心怀不平 1-12}
神實在恩待以色列那些清心的人。
2 至於我，我的腳幾乎失閃，我的腳險些滑跌。
3 我見惡人和狂傲人享平安，就心懷不平。
4 他們死的時候沒有疼痛，他們的力氣卻也壯實。
5 他們不像別人受苦，也不像別人遭災。
6 所以驕傲如鏈子戴在他們的項上，強暴像衣裳遮住他們的身體。
7 他們的眼睛因體胖而凸出，他們所得的過於心裡所想的。
8 他們譏笑人，憑惡意說欺壓人的話，他們說話自高。
9 他們的口褻瀆上天，他們的舌毀謗全地。
10 所以神的民歸到這裡，喝盡了滿杯的苦水。
11 他們說：「神怎能曉得？至高者豈有知識呢？」
12 看哪，這就是惡人！他們既是常享安逸，財寶便加增。
\subsubsection{進了神的聖所，明白了惡人的結局 13-20}
\begin{table}[H]
\centering
\begin{tabular}{p{3.75cm}|p{4.00cm}|p{9.00cm}}\hline\hline
\multicolumn{1}{c|}{抱怨徒然潔淨了心13-14}&\multicolumn{1}{c|}{難明白惡人昌盛15-16}&\multicolumn{1}{c}{在聖所看到惡人結局 17-20}\\\hline
我實在徒然潔淨了我的心，徒然洗手表明無辜！
14 因為我終日遭災難，每早晨受懲治。
&我若說我要這樣講，這就是以奸詐待你的眾子。
16 我思索怎能明白這事，眼看實係為難，
&等我進了神的聖所，思想他們的結局。
18 你實在把他們安在滑地，使他們掉在沉淪之中。
19 他們轉眼之間成了何等的荒涼，他們被驚恐滅盡了。
20 人睡醒了怎樣看夢，主啊，你醒了也必照樣輕看他們的影像。\\\hline
\end{tabular}
\end{table} 

\subsubsection{讲述看到惡人結局後,受到安慰，自述要親近神  21-28}
因而我心裡發酸，肺腑被刺。
22 我這樣愚昧無知，在你面前如畜類一般。
23 然而，我常與你同在，你攙著我的右手。
24 你要以你的訓言引導我，以後必接我到榮耀裡。
25 除你以外，在天上我有誰呢？除你以外，在地上我也沒有所愛慕的。
26 我的肉體和我的心腸衰殘，但神是我心裡的力量，又是我的福分，直到永遠。
27 遠離你的必要死亡，凡離棄你行邪淫的你都滅絕了。
28 但我親近神是與我有益！我以主耶和華為我的避難所，好叫我述說你一切的作為。


\subsection{第74篇(亞薩)：讲述仇敌拆毁圣所的恶行，求神拯救  74:1-23}
\subsubsection{求神眷顾日久荒涼之地，观看仇敵在聖所中所行的一切惡事 1-10}
神啊，你為何永遠丟棄我們呢？你為何向你草場的羊發怒，如煙冒出呢？
2 求你記念你古時所得來的會眾，就是你所贖做你產業支派的，並記念你向來所居住的錫安山。
3 求你舉步去看那日久荒涼之地，仇敵在聖所中所行的一切惡事。
4 你的敵人在你會中吼叫，他們豎了自己的旗為記號。
5 他們好像人揚起斧子，砍伐林中的樹。
6 聖所中一切雕刻的，他們現在用斧子錘子打壞了。
7 他們用火焚燒你的聖所，褻瀆你名的居所，拆毀到地。
8 他們心裡說：「我們要盡行毀滅！」他們就在遍地把神的會所都燒毀了。
9 我們不見我們的標幟，不再有先知，我們內中也沒有人知道，這災禍要到幾時呢。
10 神啊，敵人辱罵要到幾時呢？仇敵褻瀆你的名要到永遠嗎？
\subsubsection{讲述神过去的大能和作为  11-17}
你為什麼縮回你的右手？求你從懷中伸出來毀滅他們！
12 神自古以來為我的王，在地上施行拯救。
13 你曾用能力將海分開，將水中大魚的頭打破。
14 你曾砸碎鱷魚的頭，把牠給曠野的禽獸 a為食物。
15 你曾分裂磐石，水便成了溪河；你使長流的江河乾了。
16 白晝屬你，黑夜也屬你，亮光和日頭是你所預備的。
17 地的一切疆界是你所立的，夏天和冬天是你所定的。
\subsubsection{向神的祈祷,求神記念仇敵辱罵  18-23}
耶和華啊，仇敵辱罵，愚頑民褻瀆了你的名，求你記念這事。
19 不要將你斑鳩的性命交給野獸，不要永遠忘記你困苦人的性命。
20 求你顧念所立的約，因為地上黑暗之處都滿了強暴的居所。
21 不要叫受欺壓的人蒙羞回去，要叫困苦窮乏的人讚美你的名。
22 神啊，求你起來為自己申訴，要記念愚頑人怎樣終日辱罵你。
23 不要忘記你敵人的聲音，那起來敵你之人的喧嘩時常上升。

\subsection{第75篇(亞薩)：要赞美传扬神的作为，因他必按照公义审判  75:1-10}
\begin{table}[H]
\centering
\begin{tabular}{p{2cm}|p{12.00cm}|p{3.00cm}}\hline\hline
\multicolumn{1}{c|}{会众稱謝1}&\multicolumn{1}{c|}{神必按正直施行審判2-8}&\multicolumn{1}{c}{ 亞薩要歌颂神9-10}\\\hline
\multirow{2}{2.2cm}{神啊，我們稱謝你，我們稱謝你！因為你的名相近，人都述說你奇妙的作為。}
&「我到了所定的日期，必按正直施行審判。
3 地和其上的居民都消化了，我曾立了地的柱子。（細拉）
4 我對狂傲人說：『不要行事狂傲！』對凶惡人說：『不要舉角，
5 不要把你們的角高舉，不要挺著頸項說話！』」
&\multirow{2}{3.2cm}{但我要宣揚，直到永遠，我要歌頌雅各的神。
10 「惡人一切的角，我要砍斷；唯有義人的角，必被高舉。」}\\\cline{2-2}
&因為高舉非從東，非從西，也非從南而來；
7 唯有神斷定，他使這人降卑，使那人升高。
8 耶和華手裡有杯，其中的酒起沫，杯內滿了摻雜的酒。他倒出來，地上的惡人必都喝這酒的渣滓，而且喝盡。
&\\\hline
\end{tabular}
\end{table} 


\subsection{第76篇(亞薩)：神的名在地上的百姓中被尊崇，神從天的審判带来人的敬畏  76:1-12}
\subsubsection{神的名在地上的百姓中被尊崇  1-7}
\begin{table}[H]\centering
\begin{tabular}{p{8.cm}|p{8.1cm}}\hline\hline
\multicolumn{1}{c|}{神的名在地上的百姓中被尊崇  1-4}&\multicolumn{1}{c}{世上最勇敢的人都要死，无人能与神相比5-7}\\\hline
在猶大，神為人所認識；在以色列，他的名為大。
在撒冷有他的帳幕，在錫安有他的居所。
&\multirow{2}{8.2cm}{心中勇敢的人都被搶奪，他們睡了長覺，沒有一個英雄能措手。
雅各的神啊，你的斥責一發，坐車的、騎馬的都沉睡了。
唯獨你是可畏的，你怒氣一發，誰能在你面前站得住呢？}\\\cline{1-1}
他在那裡折斷弓上的火箭並盾牌、刀劍和爭戰的兵器。（細拉）
你從有野食之山而來，有光華和榮美。&\\\hline
\end{tabular}
\end{table} 

\subsubsection{神從天上使世人聽判斷 8-12}
\begin{table}[H]\centering
\begin{tabular}{p{8.75cm}|p{8.7cm}}\hline\hline
\multicolumn{1}{c|}{神审判的怒气，要带来人对祂的敬畏8-10}&\multicolumn{1}{c}{要向神许愿并要敬畏神 11-12}\\\hline
你從天上使人聽判斷——
9 神起來施行審判，要救地上一切謙卑的人——那時地就懼怕而靜默。（細拉）
10 人的憤怒要成全你的榮美，人的餘怒你要禁止。
&你們許願，當向耶和華你們的神還願；在他四面的人，都當拿貢物獻給那可畏的主。
12 他要挫折王子的驕氣，他向地上的君王顯威可畏。\\\hline
\end{tabular}
\end{table} 


\subsection{第77篇(亞薩)：在迫切寻求中因默想神过去的作为而得坚定  77:1-20}
\subsubsection{患難悲傷之日迫切尋求神,怀疑神止住了怜悯和慈爱  1-9}
\begin{table}[H]\centering
\begin{tabular}{p{8.75cm}|p{8.7cm}}\hline\hline
\multicolumn{1}{c|}{患難悲傷之日迫切尋求神的心情}&\multicolumn{1}{c}{追想古時,怀疑神止住怜悯和慈爱}\\\hline
我要向神發聲呼求，我向神發聲，他必留心聽我。
我在患難之日尋求主，我在夜間不住地舉手禱告，我的心不肯受安慰。
我想念神，就煩躁不安；我沉吟悲傷，心便發昏。（細拉）
你叫我不能閉眼，我煩亂不安，甚至不能說話。
&我追想古時之日，上古之年。
我想起我夜間的歌曲，捫心自問，我心裡也仔細省察：
難道主要永遠丟棄我，不再施恩嗎？
難道他的慈愛永遠窮盡，他的應許世世廢棄嗎？
難道神忘記開恩，因發怒就止住他的慈悲嗎？（細拉）\\\hline
\end{tabular}
\end{table} 

\subsubsection{默念神过去的作為和大能時，亚萨的信心被坚定 10-20}
\begin{table}[H]
\centering
\begin{tabular}{p{6cm}|p{11.00cm}}\hline\hline
\multicolumn{1}{c|}{提醒自己要默念神古時的作为10-13}&\multicolumn{1}{c}{神古时的奇事作为14-20}\\\hline
\multirow{3}{6.1cm}{我便說：「這是我的懦弱，但我要追念至高者顯出右手之年代。」
我要提說耶和華所行的，我要記念你古時的奇事。你是行奇事的神，
我也要思想你的經營，默念你的作為。神啊，你的作為是潔淨的，有何神大如神呢？}
&你曾在列邦中彰顯你的能力。
你曾用你的膀臂贖了你的民，就是雅各和約瑟的子孫。（細拉）\\\cline{2-2}
&神啊，諸\textcolor{red}{水}見你，一見就都驚惶，深淵也都戰抖。
\textcolor{red}{雲}中倒出水來，天空發出響聲，你的\textcolor{red}{箭}也飛行四方。
你的\textcolor{red}{雷}聲在旋風中，\textcolor{red}{電光}照亮世界，大地戰抖震動。
你的道在\textcolor{red}{海}中，你的路在大水中，你的腳蹤無人知道。\\\cline{2-2}
&你曾藉摩西和亞倫的手引導你的百姓，好像羊群一般。\\\hline
\end{tabular}
\end{table} 


\subsection{第78篇(亞薩)：提醒记住并流传百姓屡次的悖逆和神的拯救作为  78:1-72}
\begin{itemize}
\itemsep-.45em  
\item 亞薩提醒百姓留心听并时代铭记流传过去的悖逆和神的拯救 1-8
\item 叙述百姓屡次的悖逆及神不停的拯救  9-72
\begin{itemize}
\itemsep-.45em  
\item 以法莲的子孙临阵退后不守神的约 9-16
\item 在干燥之地试探神，求吃肉 17-31
\item 32-39
\item 40-55
\item 56-72
\end{itemize}
\end{itemize}

\subsubsection{要留心聽神的律法和作為并要傳給子孫，不要像祖宗悖逆 1-11}
\begin{table}[H]
\hspace{-6mm}
\centering
\begin{tabular}{p{6.75cm}|p{10.75cm}}\hline\hline
\multicolumn{1}{c|}{要留心聽神的律法和作為1-4}&\multicolumn{1}{c}{要傳給子孫，不要像祖宗悖逆5-11}\\\hline\hline
我的民哪，你們要留心聽我的訓誨，側耳聽我口中的話。
我要開口說比喻，我要說出古時的謎語，
是我們所聽見所知道的，也是我們的祖宗告訴我們的。
我們不將這些事向他們的子孫隱瞞，要將耶和華的美德和他的能力，並他奇妙的作為，述說給後代聽。
&因為他在雅各中立法度，在以色列中設律法，是他吩咐我們祖宗要傳給子孫的，
使將要生的後代子孫可以曉得。他們也要起來告訴他們的子孫，
好叫他們仰望神，不忘記神的作為，唯要守他的命令，
不要像他們的祖宗，是頑梗悖逆居心不正之輩，向著神心不誠實。
以法蓮的子孫帶著兵器拿著弓，臨陣之日轉身退後。
他們不遵守神的約，不肯照他的律法行。
又忘記他所行的，和他顯給他們奇妙的作為。\\\hline
\end{tabular}
\end{table}

\subsubsection{神在曠野的奇事救贖并祖宗悖逆試探神 12-41}
\begin{table}[H]
\hspace{-6mm}
\centering
\begin{tabular}{p{2.75cm}|p{8.25cm}|p{5.75cm}}\hline\hline
\multicolumn{1}{c|}{出埃的奇事12-16}&\multicolumn{1}{c|}{因食物試探神17-31}&\multicolumn{1}{c}{不信又說謊再三試探神32-41}\\\hline\hline
他在埃及地，在瑣安田，在他們祖宗的眼前，施行奇事。
13 他將海分裂，使他們過去，又叫水立起如壘。
14 他白日用雲彩，終夜用火光，引導他們。
他在曠野分裂磐石，多多地給他們水喝，如從深淵而出。
16 他使水從磐石湧出，叫水如江河下流。
&\multirow{2}{8.4cm}{他們卻仍舊得罪他，在乾燥之地悖逆至高者。
他們心中試探神，隨自己所欲的求食物，
並且妄論神說：「神在曠野豈能擺設筵席嗎？
他曾擊打磐石，使水湧出成了江河，他還能賜糧食嗎？還能為他的百姓預備肉嗎？」
所以耶和華聽見就發怒，有烈火向雅各燒起，有怒氣向以色列上騰。
因為他們不信服神，不倚賴他的救恩。
他卻吩咐天空，又敞開天上的門，
\textcolor{red}{降嗎哪}像雨給他們吃，將天上的糧食賜給他們。
各人(人)吃大能者的食物，他賜下糧食，使他們飽足。
他領東風起在天空，又用能力引了南風來。
他\textcolor{red}{降肉}像雨在他們當中，多如塵土，又降飛鳥，多如海沙，
落在他們的營中，在他們住處的四面。
他們吃了，而且飽足，這樣就隨了他們所欲的。
他們貪而無厭，食物還在他們口中的時候，
神的怒氣就向他們上騰，殺了他們內中的肥壯人，打倒以色列的少年人。}
&\multirow{2}{5.9cm}{雖是這樣，他們仍舊犯罪，不信他奇妙的作為。
因此他叫他們的日子全歸虛空，叫他們的年歲盡屬驚恐。
他殺他們的時候，他們才求問他，回心轉意，切切地尋求神。
他們也追念神是他們的磐石，至高的神是他們的救贖主。
他們卻用口諂媚他，用舌向他說謊。
因他們的心向他不正，在他的約上也不忠心。
但他有憐憫，赦免他們的罪孽，不滅絕他們，而且屢次消他的怒氣，不發盡他的憤怒。
他想到他們不過是血氣，是一陣去而不返的風。
他們在曠野悖逆他，在荒地叫他擔憂，何其多呢！
他們再三試探神，惹動以色列的聖者。}\\\hline
\end{tabular}
\end{table}

\subsubsection{神在埃及顯神蹟,領他們穩妥渡海 42-53}
\begin{table}[H]
\hspace{-6mm}
\centering
\begin{tabular}{p{14.5cm}|p{3.cm}}\hline\hline
\multicolumn{1}{c|}{埃及地顯十災神蹟42-51}&\multicolumn{1}{c}{領民如羊52-53}\\\hline\hline
他們不追念他的能力(手)和贖他們脫離敵人的日子。
他怎樣在埃及地顯神蹟，在瑣安田顯奇事，
把他們的江河並河汊的\textcolor{red}{水都變為血}，使他們不能喝。
他叫\textcolor{red}{蒼蠅}成群落在他們當中，嘬盡他們，又叫\textcolor{red}{青蛙}滅了他們；
把他們的土產交給螞蚱，把他們辛苦得來的交給\textcolor{red}{蝗蟲}。
他降\textcolor{red}{冰雹}打壞他們的葡萄樹，下嚴霜打壞他們的桑樹；
又把他們的牲畜交給冰雹，把他們的群畜交給閃電。
他使猛烈的怒氣和憤怒、惱恨、苦難成了一群降災的使者，臨到他們。
他為自己的怒氣修平了路，將他們交給\textcolor{red}{瘟疫}，使他們死亡。
在埃及\textcolor{red}{擊殺一切長子}，在含的帳篷中擊殺他們強壯時頭生的。
&他卻領出自己的民如羊，在曠野引他們如羊群。
他領他們穩穩妥妥的，使他們不致害怕，海卻淹沒他們的仇敵。\\\hline
\end{tabular}
\end{table}


\subsubsection{神在曠野的奇事救贖并祖宗悖逆試探神 54-72}
\begin{table}[H]
\hspace{-6mm}
\centering
\begin{tabular}{p{3.cm}|p{9.25cm}|p{5cm}}\hline\hline
\multicolumn{1}{c|}{54-55}&\multicolumn{1}{c|}{56-66}&\multicolumn{1}{c}{67-72}\\\hline\hline
他帶他們到自己聖地的邊界，到他右手所得的這山地。
55 他在他們面前趕出外邦人，用繩子將外邦的地量給他們為業，叫以色列支派的人住在他們的帳篷裡。
&他們仍舊試探、悖逆至高的神，不守他的法度，
反倒退後，行詭詐，像他們的祖宗一樣，他們改變如同翻背的弓。
因他們的丘壇惹了他的怒氣，因他們雕刻的偶像觸動他的憤恨。
神聽見就發怒，極其憎惡以色列人。
甚至他離棄示羅的帳幕，就是他在人間所搭的帳篷；
又將他的約櫃(能力)交於人擄去，將他的榮耀交在敵人手中。
並將他的百姓交於刀劍，向他的產業發怒。
少年人被火燒滅，處女也無喜歌。
祭司倒在刀下，寡婦卻不哀哭。
那時主像世人睡醒，像勇士飲酒呼喊。
他就打退了他的敵人，叫他們永蒙羞辱。
&並且他棄掉約瑟的帳篷，不揀選以法蓮支派，
卻揀選猶大支派，他所喜愛的錫安山，
蓋造他的聖所好像高峰，又像他建立永存之地。
又揀選他的僕人大衛，從羊圈中將他召來，
叫他不再跟從那些帶奶的母羊，為要牧養自己的百姓雅各和自己的產業以色列。
於是他按心中的純正牧養他們，用手中的巧妙引導他們。\\\hline
\end{tabular}
\end{table}


\subsection{第79篇(亞薩)：哀訴聖殿荒凉，求神拯救使外邦知道神，百姓讚美神  79:1-13}


\subsubsection{哀訴耶路撒冷屠毀之苦，问神幾時结束对百姓的愤怒 1-5}
神啊，外邦人進入你的產業，汙穢你的聖殿，使耶路撒冷變成荒堆，
2 把你僕人的屍首交於天空的飛鳥為食，把你聖民的肉交於地上的野獸，
3 在耶路撒冷周圍流他們的血如水，無人葬埋。
4 我們成為鄰國的羞辱，成為我們四圍人的嗤笑譏刺。
5 耶和華啊，這到幾時呢？你要動怒到永遠嗎？你的憤恨要如火焚燒嗎？

\subsubsection{求神拯救赦免，使外邦知道神，百姓讚美神 6-13}
\begin{table}[H]
\centering
\begin{tabular}{p{3.25cm}|p{5cm}|p{8.25cm}}\hline\hline
\multicolumn{1}{c|}{将憤怒倒在外邦6-7}&\multicolumn{1}{c|}{拯救赦免百姓8-9}&\multicolumn{1}{c}{使外邦知道神，百姓称谢神10-13}\\\hline
願你將你的憤怒倒在那不認識你的外邦，和那不求告你名的國度，
因為他們吞了雅各，把他的住處變為荒場。
&求你不要記念我們先祖的罪孽，向我們追討。願你的慈悲快迎著我們，因為我們落到極卑微的地步！
拯救我們的神啊，求你因你名的榮耀幫助我們，為你名的緣故搭救我們，赦免我們的罪。
&為何容外邦人說：「他們的神在哪裡呢？」願你使外邦人知道你在我們眼前，申你僕人流血的冤。
願被囚之人的嘆息達到你面前，願你按你的大能力存留那些將要死的人。
主啊，願你將我們鄰邦所羞辱你的羞辱加七倍歸到他們身上。
這樣，你的民，你草場的羊，要稱謝你直到永遠，要述說讚美你的話直到萬代。\\\hline
\end{tabular}
\end{table} 


\subsection{第80篇(亞薩)：神從前和现在对百姓的态度，求神回轉拯救,使百姓回轉得救 1-19}
\subsubsection{神现在对百姓發怒 1-6}
領約瑟如領羊群之以色列的牧者啊，求你留心聽！坐在二基路伯上的啊，求你發出光來！
2 在以法蓮、便雅憫、瑪拿西前面施展你的大能，來救我們！
3 神啊，\textcolor{red}{求你使我們回轉（復興）}；使你的臉發光，我們便要得救！
4 耶和華萬軍之神啊，你向你百姓的禱告發怒，要到幾時呢？
5 你以眼淚當食物給他們吃，又多量出眼淚給他們喝。
6 你使鄰邦因我們紛爭，我們的仇敵彼此戲笑。

\subsubsection{神從埃及挪出一棵葡萄樹 7-11}
7 萬軍之神啊，\textcolor{red}{求你使我們回轉}；使你的臉發光，我們便要得救！
8 你從埃及挪出一棵葡萄樹，趕出外邦人，把這樹栽上。
9 你在這樹跟前預備了地方，他就深深扎根，爬滿了地。
10 它的影子遮滿了山，枝子好像佳美的香柏樹。
11 它發出枝子長到大海，發出蔓子延到大河。

\subsubsection{求神回轉眷顧保護這棵葡萄樹 12-15}
12 你為何拆毀這樹的籬笆，任憑一切過路的人摘取？
13 林中出來的野豬把他糟蹋，野地的走獸拿他當食物。
14 萬軍之神啊，\textcolor{red}{求你回轉}，從天上垂看，眷顧這葡萄樹，
15 保護你右手所栽的和你為自己所堅固的枝子。

\subsubsection{求神拯救使百姓回轉 16-19}
16 這樹已經被火焚燒，被刀砍伐，他們因你臉上的怒容就滅亡了。
17 願你的手扶持你右邊的人，就是你為自己所堅固的人子。
18 這樣，我們便不退後離開你。求你救活我們，我們就要求告你的名。
19 耶和華萬軍之神啊，\textcolor{red}{求你使我們回轉}；使你的臉發光，我們便要得救！


\subsection{第81篇(亞薩)： 呼吁百姓赞美神，重申神对百姓說的话  81:1-16}
\subsubsection{呼吁百姓赞美神，并纪念神设立的律例节期 1-5}
亞薩的詩，交於伶長。用迦特樂器。
1 你們當向神我們的力量大聲歡呼，向雅各的神發聲歡樂！
2 唱起詩歌，打手鼓，彈美琴與瑟！
3 當在月朔並月望，我們過節的日期吹角！
4 因這是為以色列定的律例，是雅各神的典章。
5 他去攻擊埃及地的時候，在約瑟中間立此為證。我在那裡聽見我所不明白的言語，

\subsubsection{叙述神向百姓說的话 6-16}
\begin{table}[H]
\centering
\begin{tabular}{p{8.3cm}|p{8.5cm}}\hline\hline
\multicolumn{1}{c|}{神做的}&\multicolumn{1}{c}{民做的（悖逆和顺服）}\\\hline
\multirow{2}{8.5cm}{神說：「我使你的肩得脫重擔，你的手放下筐子。
你在急難中呼求，我就搭救你。我在雷的隱密處應允你，在米利巴水那裡試驗你。（細拉）
我的民哪，你當聽，我要勸誡你！以色列啊，甚願你肯聽從我！
在你當中不可有別的神，外邦的神你也不可下拜。
我是耶和華你的神，曾把你從埃及地領上來。你要大大張口，我就給你充滿。}
&無奈我的民不聽我的聲音，以色列全不理我。
12 我便任憑他們心裡剛硬，隨自己的計謀而行。\\\cline{2-2}
&甚願我的民肯聽從我，以色列肯行我的道！
14 我便速速制伏他們的仇敵，反手攻擊他們的敵人。
15 恨耶和華的人必來投降，但他的百姓必永久長存。
16 他也必拿上好的麥子給他們吃，又拿從磐石出的蜂蜜叫他們飽足。\\\hline
\end{tabular}
\end{table} 

\subsection{第82篇(亞薩)： 神在掌权者(至高者的兒子)中施行审判  82:1-8}
\begin{table}[H]
\centering
\begin{tabular}{p{15.25cm}|p{2.2cm}}\hline\hline
\multicolumn{1}{c|}{審判至高者的兒子}&\multicolumn{1}{c}{審判世界}\\\hline
神站在有權力者的會中，在諸神中行審判，
2 說：「你們審判不秉公義，徇惡人的情面，要到幾時呢？（細拉）
3 你們當為貧寒的人和孤兒申冤，當為困苦和窮乏的人施行公義。
4 當保護貧寒和窮乏的人，救他們脫離惡人的手。
5 你們仍不知道，也不明白，在黑暗中走來走去，地的根基都搖動了。
&\multirow{2}{2.2cm}{神啊，求你起來，審判世界，因為你要得萬邦為業}\\\cline{1-1}
我曾說：『你們是神，都是至高者的兒子。』
然而你們要死，與世人一樣；要仆倒，像王子中的一位。」&\\\hline
\end{tabular}
\end{table} 



\subsection{第83篇(亞薩)：向神控诉仇敌的恶行，求神报应敌人的祷告 83:1-18}
\subsubsection{求神除灭要剪滅以色列的仇敵 1-11}
\begin{table}[H]
\centering
\begin{tabular}{p{8.5cm}|p{4.2500cm}|p{4.50cm}}\hline\hline
\multicolumn{1}{c|}{仇敵恶行1-5}&\multicolumn{1}{c|}{仇敵名称6-8}&\multicolumn{1}{c}{求神剪滅仇敵9-11}\\\hline
神啊，求你不要靜默！神啊，求你不要閉口，也不要不作聲！
&\multirow{2}{4.4cm}{就是住帳篷的\textcolor{red}{以東人}和\textcolor{red}{以實瑪利人}，\textcolor{red}{摩押}和\textcolor{red}{夏甲人}，
\textcolor{red}{迦巴勒}、\textcolor{red}{亞捫}和\textcolor{red}{亞瑪力}、\textcolor{red}{非利士}，並\textcolor{red}{推羅}的居民。
\textcolor{red}{亞述}也與他們聯合，他們做羅得子孫的幫手。（細拉）}
&\multirow{2}{4.4cm}{求你待他們如\textcolor{red}{待米甸}，如在基順河\textcolor{red}{待西西拉和耶賓}一樣，
他們在隱多珥滅亡，成了地上的糞土。
求你叫他們的首領像俄立和西伊伯，叫他們的王子都像西巴和撒慕拿，}\\\cline{1-1}
因為你的仇敵喧嚷，恨你的抬起頭來。
他們同謀奸詐，要害你的百姓；彼此商議，要害你所隱藏的人。
他們說：「來吧，我們將他們剪滅，使他們不再成國，使以色列的名不再被人記念。」
他們同心商議，彼此結盟，要抵擋你。&&\\\hline
\end{tabular}
\end{table} 

\subsubsection{求神除灭要得神的住處的仇敵 1-11}
\begin{table}[H]
\centering
\begin{tabular}{p{2.5cm}|p{7.300cm}|p{7.7cm}}\hline\hline
\multicolumn{1}{c|}{仇敵}&\multicolumn{1}{c|}{求神除灭}&\multicolumn{1}{c}{願你使他們}\\\hline
他們說：「我們要得神的住處，作為自己的產業。」
&我的神啊，求你叫他們像旋風的塵土，像風前的碎秸。
火怎樣焚燒樹林，火焰怎樣燒著山嶺，
求你也照樣用狂風追趕他們，用暴雨恐嚇他們。
&願你使他們滿面羞恥好叫他們尋求你耶和華的名。
願他們永遠羞愧驚惶，願他們慚愧滅亡，
使他們知道，唯獨你名為耶和華的是全地以上的至高者。\\\hline
\end{tabular}
\end{table} 


\subsection{第84篇(可拉後裔)： 住在神殿,向往锡安大道,倚靠神的有福  84:1-12}
\begin{table}[H]
\centering
\begin{tabular}{p{5.9cm}|p{3.6cm}|p{7.900cm}}\hline\hline
\multicolumn{1}{c|}{渴想住在神殿1-4}&\multicolumn{1}{c|}{行走錫安大道5-7}&\multicolumn{1}{c}{求告萬軍之神8-12}\\\hline
萬軍之耶和華啊，你的居所何等可愛！
我羨慕渴想耶和華的院宇，我的心腸、我的肉體向永生神呼籲(歡呼)。
萬軍之耶和華，我的王我的神啊，在你祭壇那裡，麻雀為自己找著房屋，燕子為自己找著抱雛之窩。
如此住在你殿中的，便為有福！他們仍要讚美你。（細拉）
&靠你有力量，心中想往錫安大道的，這人便為有福！
他們經過流淚谷，叫這谷變為泉源之地，並有秋雨之福蓋滿了全谷。
他們行走，力上加力，各人到錫安朝見神。
&耶和華萬軍之神啊，求你聽我的禱告！雅各的神啊，求你留心聽！（細拉）
神啊，你是我們的盾牌，求你垂顧觀看你受膏者的面。
在你的院宇住一日，勝似在別處住千日。寧可在我神殿中看門，不願住在惡人的帳篷裡。
因為耶和華神是日頭，是盾牌，要賜下恩惠和榮耀，他未嘗留下一樣好處，不給那些行動正直的人。
萬軍之耶和華啊，倚靠你的人便為有福！\\\hline
\end{tabular}
\end{table} 


\subsection{第85篇(可拉後裔)： 祷告求神再次施恩，并坚信神必应许祷告 85:1-13}
\subsubsection{神已向雅各施恩,求神再次收回怒气，向百姓施行慈爱 1-7}
\begin{table}[H]
\centering
\begin{tabular}{p{6.75cm}|p{10.00cm}}\hline\hline
\multicolumn{1}{c|}{神已向雅各施恩}&\multicolumn{1}{c}{求神再次收回怒气，向百姓施行慈爱}\\\hline
耶和華啊，你已經向你的地施恩，救回被擄的雅各。
2 你赦免了你百姓的罪孽，遮蓋了他們一切的過犯。（細拉）
3 你收轉了所發的憤怒，和你猛烈的怒氣。
&拯救我們的神啊，求你使我們回轉，叫你的惱恨向我們止息。你要向我們發怒到永遠嗎？你要將你的怒氣延留到萬代嗎？
6 你不再將我們救活，使你的百姓靠你歡喜嗎？
7 耶和華啊，求你使我們得見你的慈愛，又將你的救恩賜給我們。\\\hline
\end{tabular}
\end{table} 

\subsubsection{神必應許将平安賜給百姓 8-13}
8 我要聽神耶和華所說的話，因為他必應許將平安賜給他的百姓他的聖民，他們卻不可再轉去妄行。
9 他的救恩誠然與敬畏他的人相近，叫榮耀住在我們的地上。
10 慈愛和誠實彼此相遇，公義和平安彼此相親。
11 誠實從地而生，公義從天而現。
12 耶和華必將好處賜給我們，我們的地也要多出土產。
13 公義要行在他面前，叫他的腳蹤成為可走的路。

\subsection{第86篇(大卫)： 求神拯救脱离仇敌，免下阴间的祷告 86:1-17}
\subsubsection{求神在患难之日拯救的祷告 1-7}
耶和華啊，求你側耳應允我，因我是困苦窮乏的。
2 求你保存我的性命，因我是虔誠人。我的神啊，求你拯救這倚靠你的僕人。
3 主啊，求你憐憫我，因我終日求告你。
4 主啊，求你使僕人心裡歡喜，因為我的心仰望你。
5 主啊，你本為良善，樂意饒恕人，有豐盛的慈愛賜給凡求告你的人。
6 耶和華啊，求你留心聽我的禱告，垂聽我懇求的聲音。
7 我在患難之日要求告你，因為你必應允我
\subsubsection{求神引导免下阴间,萬民都要來敬拜榮耀你 8-13}
主啊，諸神之中沒有可比你的，你的作為也無可比。
9 主啊，你所造的萬民都要來敬拜你，他們也要榮耀你的名。
10 因你為大，且行奇妙的事，唯獨你是神。
11 耶和華啊，求你將你的道指教我，我要照你的真理行；求你使我專心敬畏你的名。
12 主我的神啊，我要一心稱讚你，我要榮耀你的名直到永遠。
13 因為你向我發的慈愛是大的，你救了我的靈魂免入極深的陰間
\subsubsection{求神救自己脱离仇敌，使敌人羞愧 14-17}
14 神啊，驕傲的人起來攻擊我，又有一黨強橫的人尋索我的命，他們沒有將你放在眼中。
15 主啊，你是有憐憫有恩典的神，不輕易發怒，並有豐盛的慈愛和誠實。
16 求你向我轉臉，憐恤我，將你的力量賜給僕人，救你婢女的兒子。
17 求你向我顯出恩待我的憑據，叫恨我的人看見便羞愧，因為你耶和華幫助我，安慰我。

\subsection{第87篇(可拉後裔)： 錫安榮耀的事-神要在各族点出生在锡安的人 87:1-7}
\begin{table}[H]
\centering
\begin{tabular}{p{5.25cm}|p{12.00cm}}\hline\hline
\multicolumn{1}{c|}{神爱锡安 1-3}&\multicolumn{1}{c}{錫安榮耀的事 4-7}\\\hline
耶和華所立的根基在聖山上。
他愛錫安的門，勝於愛雅各一切的住處。
神的城啊，有榮耀的事乃指著你說的。（細拉）
&我要提起\textcolor{red}{拉哈伯}和\textcolor{red}{巴比倫}人，是在認識我之中的；看哪，\textcolor{red}{非利士}和\textcolor{red}{推羅}並\textcolor{red}{古實}人，個個生在那裡。
論到錫安，必說：這一個、那一個都生在其中，而且至高者必親自堅立這城。
當耶和華記錄萬民的時候，他要點出這一個生在那裡。（細拉）
歌唱的，跳舞的，都要說：我的泉源都在你裡面。\\\hline
\end{tabular}
\end{table} 

\subsection{第88篇（可拉後裔：以斯拉人希幔）：因神的憤怒心中愁苦，而迫切天天呼求神  88:1-18}
\subsubsection{因神的憤怒身处患難，而迫切天天呼求神 1-9}
\begin{table}[H]
\centering
\begin{tabular}{p{4.75cm}|p{12.00cm}}\hline\hline
\multicolumn{1}{c|}{禱告神，求神聽呼求}&\multicolumn{1}{c}{因神的憤怒，使得自己愁苦}\\\hline
可拉後裔的詩歌，就是以斯拉人希幔的訓誨詩，交於伶長。調用麻哈拉利暗俄。
1 耶和華拯救我的神啊，我晝夜在你面前呼籲。
2 願我的禱告達到你面前，求你側耳聽我的呼求。
&因為我心裡滿了患難，我的性命臨近陰間。
4 我算和下坑的人同列，如同無力 (沒有幫助)的人一樣。
5 我被丟在死人中，好像被殺的人躺在墳墓裡，他們是你不再記念的，與你隔絕了。
6 你把我放在極深的坑裡，在黑暗地方，在深處。
7 你的憤怒重壓我身，你用一切的波浪困住我。（細拉）
8 你把我所認識的隔在遠處，使我為他們所憎惡。我被拘困，不得出來。
9 我的眼睛因困苦而乾癟。耶和華啊，我天天求告你，向你舉手。\\\hline
\end{tabular}
\end{table} 

\subsubsection{盼神早日垂听祷告，再述自己处境艰难 10-18}
\begin{table}[H]
\centering
\begin{tabular}{p{10.25cm}|p{7.00cm}}\hline\hline
\multicolumn{1}{c|}{盼神早日垂听祷告}&\multicolumn{1}{c}{再述自己处境艰难}\\\hline
你豈要行奇事給死人看嗎？難道陰魂還能起來稱讚你嗎？（細拉）
11 豈能在墳墓裡述說你的慈愛嗎？豈能在滅亡中述說你的信實嗎？
12 你的奇事豈能在幽暗裡被知道嗎？你的公義豈能在忘記之地被知道嗎？
13 耶和華啊，我呼求你，我早晨的禱告要達到你面前。
14 耶和華啊，你為何丟棄我？為何掩面不顧我？
&我自幼受苦，幾乎死亡；我受你的驚恐，甚至慌張。
16 你的烈怒漫過我身，你的驚嚇把我剪除。
17 這些終日如水環繞我，一齊都來圍困我。
18 你把我的良朋密友隔在遠處，使我所認識的人進入黑暗裡。\\\hline
\end{tabular}
\end{table} 


\subsection{第89篇（以斯拉人以探)：神尋得大衛并与之立基督国度的約，人要歌唱神的信实 1-52}
\subsubsection{神拣选大卫并与之立约，使人更明白神的属性而讚美神 30-52}
\begin{table}[H]
\centering
\begin{tabular}{p{3.cm}|p{7.75cm}|p{6.00cm}}\hline\hline
\multicolumn{1}{c|}{稱讚神立大衛之約}&\multicolumn{1}{c}{歌唱神的属性}&\multicolumn{1}{c}{立大衛之約的过程}\\\hline
我要歌唱耶和華的慈愛直到永遠，我要用口將你的信實傳於萬代。
2 因我曾說：「你的慈悲必建立到永遠，你的信實必堅立在天上。」

3 「我與我所揀選的人立了約，向我的僕人大衛起了誓：
4 我要建立你的後裔直到永遠，要建立你的寶座直到萬代。」（細拉）
5 耶和華啊，諸天要稱讚你的奇事，在聖者的會中要稱讚你的信實！
&在天空，誰能比耶和華呢？神的眾子中，誰能像耶和華呢？
7 他在聖者的會中是大有威嚴的神，比一切在他四圍的更可畏懼。
8 耶和華萬軍之神啊，哪一個大能者像你耶和華？你的信實是在你的四圍。

9 你管轄海的狂傲，波浪翻騰，你就使他平靜了。
10 你打碎了拉哈伯，似乎是已殺的人，你用有能的膀臂打散了你的仇敵。

11 天屬你，地也屬你，世界和其中所充滿的都為你所建立。
12 南北為你所創造，他泊和黑門都因你的名歡呼。
13 你有大能的膀臂，你的手有力，你的右手也高舉。

公義和公平是你寶座的根基，慈愛和誠實行在你前面。
15 知道向你歡呼的，那民是有福的！耶和華啊，他們在你臉上的光裡行走。
16 他們因你的名終日歡樂，因你的公義得以高舉。
17 你是他們力量的榮耀，因為你喜悅我們，我們的角必被高舉。
18 我們的盾牌屬耶和華，我們的王屬以色列的聖者。
& 當時，你在異象中曉諭你的聖民說：「我已把救助之力加在那有能者的身上，我高舉那從民中所揀選的。
20 我尋得我的僕人大衛，用我的聖膏膏他。
21 我的手必使他堅立，我的膀臂也必堅固他。
22 仇敵必不勒索他，凶惡之子也不苦害他。
23 我要在他面前打碎他的敵人，擊殺那恨他的人。
24 只是我的信實和我的慈愛要與他同在，因我的名他的角必被高舉。
25 我要使他的左手伸到海上，右手伸到河上。
26 他要稱呼我說：『你是我的父，是我的神，是拯救我的磐石。』
27 我也要立他為長子，為世上最高的君王。
28 我要為他存留我的慈愛直到永遠，我與他立的約必要堅定。
29 我也要使他的後裔存到永遠，使他的寶座如天之久。\\\hline
\end{tabular}
\end{table} 


\subsubsection{对背约者的管教 30-52}
\begin{table}[H]
\centering
\begin{tabular}{p{5.5cm}|p{5.cm}|p{6.75cm}}\hline\hline
\multicolumn{1}{c|}{人背约神不背约的应许}&\multicolumn{1}{c}{主向多馬顯現}&\multicolumn{1}{c}{主向多馬顯現}\\\hline
 倘若他的子孫離棄我的律法，不照我的典章行，
31 背棄我的律例，不遵守我的誡命，
32 我就要用杖責罰他們的過犯，用鞭責罰他們的罪孽。
33 只是我必不將我的慈愛全然收回，也必不叫我的信實廢棄。
34 我必不背棄我的約，也不改變我口中所出的。
35 我一次指著自己的聖潔起誓，我決不向大衛說謊。
36 他的後裔要存到永遠，他的寶座在我面前如日之恆一般。
37 又如月亮永遠堅立，如天上確實的見證。」（細拉）
&但你惱怒你的受膏者，就丟掉棄絕他。
39 你厭惡了與僕人所立的約，將他的冠冕踐踏於地。
40 你拆毀了他一切的籬笆，使他的保障變為荒場。
41 凡過路的人都搶奪他，他成為鄰邦的羞辱。
42 你高舉了他敵人的右手，你叫他一切的仇敵歡喜。
43 你叫他的刀劍捲刃，叫他在爭戰之中站立不住。
44 你使他的光輝止息，將他的寶座推倒於地。
45 你減少他青年的日子，又使他蒙羞。（細拉）
&耶和華啊，這要到幾時呢？你要將自己隱藏到永遠嗎？你的憤怒如火焚燒要到幾時呢？
47 求你想念，我的時候是何等地短少！你創造世人，要使他們歸何等的虛空呢！（細拉）
48 誰能常活免死，救他的靈魂脫離陰間的權柄呢？（細拉）
49 主啊，你從前憑你的信實向大衛立誓要施行的慈愛，在哪裡呢？
50 主啊，求你記念僕人們所受的羞辱，記念我怎樣將一切強盛民的羞辱存在我懷裡。
51 耶和華啊，你的仇敵用這羞辱羞辱了你的僕人，羞辱了你受膏者的腳蹤。
52 耶和華是應當稱頌的，直到永遠！阿們，阿們。\\\hline
\end{tabular}
\end{table} 
